\section{Introduction}
Self-evolving agents change the systems through which they act: they retain memories, create tools, restructure workflows, and update models. These changes can improve performance while introducing unsafe behavior. Misevolution studies this risk across these four pathways~\cite{shao2025misevolve}. Its released project illustrates a tool created to share public posters being reused for a confidential financial report, producing a public link~\cite{misevolutionrepo}.

Executable contracts make safety requirements explicit. RvLLM provides a language for domain-specific runtime specifications~\cite{zhang2025rvllm}; Agent Behavioral Contracts studies explicit behavioral constraints and enforcement~\cite{bhardwaj2026abc}. However, compliance with an incomplete contract does not establish safety. A rule requiring authorized recipients for direct delivery may remain satisfied when a newly created link exposes the same file to others. Updating the monitor's bindings alone does not add the missing access restriction.

We study \emph{policy-guided contract co-evolution}. A developer supplies high-level safety requirements \(P\) before evolution, together with an initial executable contract \(C_0\). The policy remains fixed. After an agent update, the method asks whether its effects require additional executable obligations to enforce that policy. It does not infer the policy from the update or introduce a separate latent human-intent target.

The difficult step is identifying what is missing. New and existing capabilities can interact through data flow, control flow, and shared state; inspecting new tool signatures alone misses these interactions. Exhaustively anticipating their consequences is generally infeasible. We narrow the search using two inputs: the policy identifies consequences of interest, and the agent delta identifies where dependency exploration starts. A bounded slice includes affected old and new operations, their resource inputs, and relevant guards.

Within this slice, three tests generate candidate gaps: an effect with no directly applicable clause, a necessary policy condition left unconstrained, and an altered path that bypasses an existing restriction. The tests focus validation on potentially missing conditions. Validation seeks a feasible behavior that satisfies the old contract while violating the supplied policy, accounting for existing guards and jointly applicable clauses.

Supported findings justify additions or refinements to the contract. Each edit records its policy source, the relevant change and dependency path, and validation evidence. Existing obligations are retained and checked for unintended weakening. The candidate is then checked against the updated contract; a controller repairs, restricts, or rejects unsuccessful updates. Prevention requires interception before the harmful effect.

We adopt RvLLM's Expert Specification Language (ESL): predicates have natural-language descriptions, and rules are first-order logical formulas whose conjunction forms a specification~\cite{zhang2025rvllm}. Our contribution concerns deriving and maintaining obligations, rather than defining a new language.

The approach complements SEVerA's synthesis and optimization under supplied hard contracts~\cite{banerjee2026severa}. We ask when the executable obligations themselves need extension. It also builds on the broader problem of requirements evolution~\cite{ali2011requirements}; the proposed focus is change-guided analysis and executable clause maintenance for evolving agents.

Our central research question is empirical: \emph{does policy-guided analysis of agent changes identify missing obligations that help prevent misbehavior while preserving useful evolution?} The method defines how to search and justify edits; experiments assess how effectively that search covers relevant interactions. Exhaustive anticipation of capability combinations is not a prerequisite for studying its practical benefit.

This proposal develops three elements: (1) change-guided analysis of policy-relevant effects and dependencies; (2) validated additions and refinements to ESL contracts, with preservation and deployment checks; and (3) an evaluation on the existing Misevolution suite that measures useful gap discovery, prevention beyond fixed contracts with updated instrumentation, and utility and cost. Implementation and empirical results remain to be completed.
